\section{A nonzero bounded transport satisfying the flat-only plateau constraints}
\label{sec:outside-only-cutoff-counterexample}

\paragraph{Status and scope.}
The complete post-writing author reading and two independent
whole-proof readings have passed.
The scope reservation in MA-PI-010 can be decided by an actual
counterexample.  Outside-plateau orthogonality alone does not
force a bounded equivariant map to vanish.  The inside plateau
on the coordinate-flat core is vacuous, so the same example
satisfies both plateau implications on that core.  It does
not satisfy the full cutoff equation.  No zero-location
assumption or zero-counting theorem is used.

Throughout this section \(S'=S=\{\infty\}\) over
\(\mathbb Q\), in the principal compact sector.  The global
outside-\(S\) factors are the spherical
\(\mathbf1_{\mathbb Z_p}\); there is no extra unit-shell
transition set.  Retain the actual local chart and actions
\[
 \begin{gathered}
 Vf(x)=2e^{x/2}f(e^x),\qquad
 D_yf(t)=e^{-y/2}f(e^{-y}t),\qquad VD_y=\tau_yV,\\
 \tau_yh(x)=h(x-y),\qquad
 \widehat f(\xi)=\int_{\mathbb R}f(t)e^{-2\pi it\xi}\,dt,\\
 \mathcal Mh(z)=\int_{\mathbb R}h(x)e^{zx}\,dx.
 \end{gathered}
 \tag{MA-PI.250}
\]
Let \(\mathscr I\subset\mathcal E\) be the actual centered
rational smooth theta ideal.  The Hilbert space \(\mathscr H\)
may be any one of
\[
 \begin{aligned}
 H^{\rm one}_{b,s}
   &=L^2(\mathbb R,e^{2bx}p_s(x)\,dx),
     &&-\tfrac12<b<\tfrac12,\quad s\ge0,\\
 H^{\rm sym}_{\alpha,s}
   &=L^2(\mathbb R,e^{2\alpha|x|}p_s(x)\,dx),
     &&0<\alpha<\tfrac12,\quad s\ge0,\\
 p_s(x)&=(1+x^2)^{s/2},\qquad
 J_{\mathscr H}=\overline{\mathscr I}^{\mathscr H}.
 \end{aligned}
 \tag{MA-PI.251}
\]
The first family is the exact printed realization with
\(b=A-1/2\), hence \(0<A<1\), after the unitary already
specified in MA-PI-002/007.  The second is the separately
specified corrected symmetric metric.  Each orthogonal
projector below is formed in its own displayed norm.

Let \(P_{Y,\mathscr H}\) be the orthogonal projector onto
the actual one-place global PW range at level \(Y>0\).
For any smooth profile allowed in MA-PI.212, set
\[
 \begin{gathered}
 0\le\rho_Y\le1,\qquad
 \rho_Y=1\text{ on }[0,Y/2],\qquad
 \rho_Y=0\text{ on }[2Y,\infty),\\
 B_Y=\mathcal F^{-1}M_{\rho_Y(|\xi|)}\mathcal F.
 \end{gathered}
 \tag{MA-PI.252}
\]
As before, \(\mathscr X_0\) is the even Schwartz space
with \(f(0)=\int f=0\), and
\(\mathscr X_{\rm fl}=\mathscr X_0\cap\mathcal S_{00}(\mathbb R)\).

\paragraph{MA-PI-012 (a genuine outside-plateau countermap in every stated metric).}
There is one fixed smooth convolution kernel \(k\), independent
of \(Y,\rho_Y,b,\alpha,s\), such that \(R=C_k\) extends to
a bounded injective translation-equivariant map
\[
                 R:\mathscr H\longrightarrow J_{\mathscr H}
 \tag{MA-PI.253}
\]
for every space in (MA-PI.251).  For every \(Y>0\), every
profile in (MA-PI.252), and every even Schwartz \(f\),
\[
       B_Yf=0\quad\Longrightarrow\quad P_{Y,\mathscr H}RVf=0.
 \tag{MA-PI.254}
\]
In particular, \(R\ne0\) satisfies both plateau implications
on \(\mathscr X_{\rm fl}\) for every such profile and level:
\[
 \begin{array}{ll}
 P_{Y,\mathscr H}RVf=RVf,& B_Yf=f,\quad f\in\mathscr X_{\rm fl},\\
 P_{Y,\mathscr H}RVf=0,& B_Yf=0,\quad f\in\mathscr X_{\rm fl}.
 \end{array}
 \tag{MA-PI.255}
\]
The first line has only the zero input.  At every fixed
smooth profile and level, however, the full equation
\(P_{Y,\mathscr H}RVf=RV(B_Yf)\) fails for some
\(f\in\mathscr X_{\rm fl}\).  Thus this example respects
the strict difference between flat-core plateau implications
and the full equation in MA-PI-010/011.

\begin{proof}
\emph{1. A genuine polar-free theta kernel with an upper support edge.}
Choose a nonzero real \(h\in C_c^\infty\) supported in
\((-\log2-2,-\log2-1)\).  With \(D=d/dx\), put
\[
 g=(D+\tfrac12)h,\qquad
 G(z)=\mathcal Mg(z)=(\tfrac12-z)\mathcal Mh(z),
 \qquad G(\tfrac12)=0.
 \tag{MA-PI.256}
\]
The compactly supported solution of \((D+1/2)h=0\) is
zero, so \(g\ne0\).  Its actual additive source is
\[
 f_g(t)=\tfrac12|t|^{-1/2}g(\log|t|)\quad(t\ne0),
 \qquad f_g(0)=0.
 \tag{MA-PI.257}
\]
This is an even annular Schwartz function.  By the exact
plus-half moment calculation (MA-PI.218),
\(\int_{\mathbb R}f_g=G(1/2)=0\).  Its value and Fourier
value at zero therefore vanish.  Its original global
centered theta function is
\[
 \begin{aligned}
 a(x)=\Theta f_g(x)
   &=2e^{x/2}\sum_{n\ge1}f_g(ne^x)
     =\sum_{n\ge1}n^{-1/2}g(x+\log n),\\
 a&\in\mathscr I\cap\mathcal E,\qquad
 \operatorname{supp}a\subset(-\infty,-\log2-1],\\
 \mathcal Ma(z)&=\zeta(\tfrac12+z)G(z).
 \end{aligned}
 \tag{MA-PI.258}
\]
The support assertion follows term by term from \(\log n\ge0\).
On the opposite tail, polar-free Poisson summation and the
Schwartz decay of \(\widehat f_g\) give every exponential
Schwartz estimate, including derivatives.  Thus the stated
membership in \(\mathcal E\) is the actual theta-image
membership, not a spectral definition.  The Mellin identity
first follows by absolute summation for \(\Re z>1/2\).
Both sides then continue to entire functions, since the
only pole on the right is canceled by (MA-PI.256).
The product is not identically zero; hence \(a\ne0\).

\emph{2. The exact Fourier chart and its nonvanishing local factor.}
For an arbitrary even Schwartz function \(f\), define
\[
       Wf(x)=e^{-x/2}\widehat f(e^{-x}),\qquad
       \Gamma_{\mathbb R}(s)=\pi^{-s/2}\Gamma(s/2).
 \tag{MA-PI.259}
\]
At \(+\infty\), \(Wf=O_f(e^{-x/2})\); at \(-\infty\)
it has arbitrarily rapid exponential decay, with the
corresponding derivative estimates.  The real even local
Fourier functional equation gives, for \(|\Re z|<1/2\),
\[
 \begin{aligned}
 \mathcal M(Wf)(z)
   &=\frac12\,\gamma(z)\mathcal M(Vf)(z),\\
 \gamma(z)&=
       \frac{\Gamma_{\mathbb R}(1/2-z)}
            {\Gamma_{\mathbb R}(1/2+z)}.
 \end{aligned}
 \tag{MA-PI.260}
\]
Here the local multiplier is holomorphic and nonzero
throughout the open strip.  To check the scalar and
orientation, the change \(t=e^{-x}\) gives
\(\mathcal M(Wf)(z)=\int_0^\infty\widehat f(t)t^{-z-1/2}\,dt\),
whereas \(\mathcal M(Vf)(z)=2\int_0^\infty f(t)t^{z-1/2}\,dt\).
Tate's even real factor is
\(\Gamma_{\mathbb R}(s)/\Gamma_{\mathbb R}(1-s)\)
in \(Z(f,s)=\rho(s)Z(\widehat f,1-s)\).
Taking \(s=1/2+z\) proves (MA-PI.260), with exactly
the displayed factor \(1/2\).  All Mellin integrals here
converge on the stated strip for the actual Schwartz inputs.

Define the holomorphic function
\[
 r(z)=\tfrac12\gamma(z)\zeta(\tfrac12+z)G(z),
                       \qquad |\Re z|<\tfrac12.
 \tag{MA-PI.261}
\]
This is defined by the local gamma quotient.  No reciprocal
of a global zeta function, including at its zeros, occurs
in this definition.

\emph{3. The full contour argument and boundedness in the actual norms.}
On every fixed closed vertical substrip of \(|\Re z|<1/2\),
\(r\) and every fixed derivative decrease faster than every
power of \(|\Im z|\), uniformly in the real part.
Indeed compact smooth \(g\) gives that property for \(G\).
The zeta factor and its derivatives have polynomial bounds
there.  One elementary justification is (MA-RS.10) with
\(N=1\): on a fixed strip \(0<\varepsilon\le\Re s
\le1-\varepsilon\),
\[
 \zeta(s)=\frac{s}{s-1}
                  -s\int_1^\infty\{u\}u^{-s-1}\,du.
 \tag{MA-PI.262}
\]
This gives a polynomial bound uniformly there, and fixed
small Cauchy circles give all required derivatives.
Stirling's formula for the quotient in (MA-PI.260) gives
polynomial bounds for that quotient and every derivative,
uniformly on each such closed substrip.  Its exponential
vertical factors cancel.  This is also the precise local
estimate already checked in (MA-MO.281)--(MA-MO.282).

Set
\[
 k(x)=\frac1{2\pi}\int_{\mathbb R}r(it)e^{-itx}\,dt.
 \tag{MA-PI.263}
\]
For every fixed \(-1/2<c<1/2\), integrate
\(r(z)e^{-zx}\) over the rectangle between the imaginary
axis and \(\Re z=c\).  It has no pole in that rectangle.
For each fixed \(x\), the two horizontal integrals tend
to zero by the uniform rapid vertical decrease just proved.
Consequently
\[
 e^{cx}k(x)=\frac1{2\pi}
              \int_{\mathbb R}r(c+it)e^{-itx}\,dt
                  \ \in\mathcal S(\mathbb R),
               \qquad -\tfrac12<c<\tfrac12.
 \tag{MA-PI.264}
\]
The same argument permits every derivative in \(x\);
alternatively use the Schwartz property on the right.
Given \(0\le\beta<1/2\), choose
\(\beta<\sigma<1/2\) and apply (MA-PI.264) with
\(c=\sigma\) on the positive half-line and \(c=-\sigma\)
on the negative half-line.  It follows, for all integers
\(j,N\ge0\), that
\[
 \int_{\mathbb R} e^{\beta|x|}(1+|x|)^N
                          |k^{(j)}(x)|\,dx<\infty.
 \tag{MA-PI.265}
\]
This is a weighted estimate for the derivatives of \(k\);
it does not assert smoothness of the nonsmooth weight
\(e^{\beta|x|}\) at zero.

Let \(R=C_k\), with
\((C_k u)(x)=\int k(y)u(x-y)\,dy\).
For \(v_s(x)=(1+x^2)^{s/4}\),
\(v_s(x+y)\le2^{s/4}v_s(x)v_s(y)\).
Weighted Young's inequality therefore gives
\[
 \begin{aligned}
 \|C_k u\|_{H^{\rm one}_{b,s}}
   &\le2^{s/4}\|e^{b\cdot}v_s k\|_{L^1}
                                         \|u\|_{H^{\rm one}_{b,s}},\\
 \|C_k u\|_{H^{\rm sym}_{\alpha,s}}
   &\le2^{s/4}\|e^{\alpha|\cdot|}v_s k\|_{L^1}
                                         \|u\|_{H^{\rm sym}_{\alpha,s}}.
 \end{aligned}
 \tag{MA-PI.266}
\]
All constants are finite by (MA-PI.265).  The single
kernel is thus bounded in every specified space and
commutes with every translation.  Its operator norms
are allowed to depend on the fixed norm parameters.

The holomorphic \(r\) is not identically zero.  On every
line \(b+i\mathbb R\) its zero set has measure zero.
Conjugation by \(e^{bx}\) makes \(R\) the ordinary
\(L^2\) Fourier multiplier \(r(b+it)\).  Therefore it
is injective on \(H^{\rm one}_{b,s}\).
For the symmetric space, use its continuous inclusion
in the central ordinary \(L^2\) space and the same
argument with \(b=0\).  This proves injectivity and,
in particular, nonzero boundedness in every stated metric.

\emph{4. Range in the actual ideal closure.}
For an even Schwartz \(f\), equations
(MA-PI.258), (MA-PI.260) and (MA-PI.261) give the actual
central Fourier identity
\[
                         RVf=a*Wf.
 \tag{MA-PI.267}
\]
All functions involved are integrable or square-integrable
on that line, so uniqueness of the ordinary Fourier
transform proves the equality of the actual functions.
Moreover \(Vf\) belongs continuously to every space in
(MA-PI.251): its left tail is \(O_f(e^{x/2})\) and its
right tail has arbitrary exponential decay.  The restrictions
\(b>-1/2\) and \(\alpha<1/2\) ensure local integrability
in the corresponding additive norm.  The images under
\(V\) of even annular tests contain every logarithmic
compact smooth function, hence are dense in these spaces.

Choose smooth compact cutoffs \(\chi_N\to1\) with
\(0\le\chi_N\le1\), and put \(w_N=\chi_NWf\).
The tails in (MA-PI.259) give convergence to \(Wf\)
in each of the weighted \(L^1\) norms with weights
\(e^{by}v_s(y)\) or \(e^{\alpha|y|}v_s(y)\).
Each compactly supported \(w_N\) gives the genuine source
\[
 F_N(t)=\int_{\mathbb R}w_N(y)D_yf_g(t)\,dy,
 \qquad
 \Theta F_N=a*w_N\in\mathscr I.
 \tag{MA-PI.268}
\]
The integral is over a compact family of dilations of
the fixed annular Schwartz source.  It is therefore
annular and Schwartz.  Its value at zero vanishes;
its integral is \(\int w_N(y)e^{y/2}\,dy\) times
\(\int f_g=0\).  Both polar conditions are retained.
Commutation of theta with normalized dilation and the
compact integral proves the displayed equality directly.
Weighted Young, now with the fixed \(a\in\mathscr H\),
shows \(a*w_N\to a*Wf\) in \(\mathscr H\).
Thus (MA-PI.267) belongs to \(J_{\mathscr H}\).
Density of \(V\mathcal S(\mathbb R)^{\rm even}\),
boundedness of \(R\), and closedness of \(J_{\mathscr H}\)
prove (MA-PI.253) on the entire Hilbert space.
No description of that closure by a selected zero divisor
is required.

\emph{5. The full one-place PW range and outside support.}
The actual one-place source PW range is generated by even
Schwartz \(F\) with both polar conditions and
\(\operatorname{supp}\widehat F\subset[-Y,Y]\), followed
by the smooth image closure and the specified Hilbert
closure.  This is Definition~7.3.2.3 and
(7.3.2.12)--(7.3.2.17), pp.~354--355, specialized to
\(S'=S=\{\infty\}\), and Definition~11.2.1.3,
(11.2.1.7)--(11.2.1.8), pp.~488--489.
It is the whole stated one-place range, not an auxiliary
subfamily.  For every such source, actual Poisson summation
gives
\[
 \Theta F(x)=e^{-x/2}
         \sum_{m\in\mathbb Z\setminus\{0\}}
                             \widehat F(me^{-x})=0
                                  \quad(x<-\log Y).
 \tag{MA-PI.269}
\]
There is no polar summand because \(F(0)=\widehat F(0)=0\).
The two closures preserve vanishing on this open half-line;
in the Hilbert norm this is a closed support condition.
Thus the full range of \(P_{Y,\mathscr H}\) is supported
in \([-\log Y,\infty)\).

If \(B_Yf=0\), then \(\widehat f=0\) on
\([-Y/2,Y/2]\), since \(\rho_Y=1\) there.  Consequently
\[
 \begin{aligned}
 \operatorname{supp}Wf
   &\subset(-\infty,-\log Y+\log2],\\
 \operatorname{supp}RVf
   &\subset(-\infty,-\log Y-1].
 \end{aligned}
 \tag{MA-PI.270}
\]
The second assertion follows from (MA-PI.258), (MA-PI.267)
and the ordinary support rule for these integrable
convolution factors.  Their supports are bounded above,
so the assertion also follows directly from the convolution
integral.  The supports in (MA-PI.269) and (MA-PI.270)
are disjoint, independently of the positive weight.
This proves their actual Hilbert orthogonality and (MA-PI.254).

If instead \(f\) is flat at zero and \(B_Yf=f\), then
\(\widehat f\) is compactly supported.  Fourier inversion
extends \(f\) to an entire function, whose full origin
jet vanishes.  Hence \(f=0\).  This proves the first
line of (MA-PI.255); the second is already proved.

\emph{6. Failure of the stronger full equation in every displayed metric.}
Fix any one allowed profile and level, and suppose, for
contradiction, that the full equation holds on
\(\mathscr X_{\rm fl}\).  The fixed-level approximation
(MA-PI.222)--(MA-PI.225) transfers it to \(\mathscr X_0\)
also in each metric (MA-PI.251).  To verify the extra
weight explicitly, the local squared error for
\(r_\varepsilon(t)=\chi(t/\varepsilon)f(t)\),
\(f\in\mathscr X_0\), satisfies
\[
 \begin{aligned}
 \|Vr_\varepsilon\|_{H^{\rm one}_{b,s}}^2
   &\le C\varepsilon^{5+2b}(1+|\log\varepsilon|)^s,\\
 \|Vr_\varepsilon\|_{H^{\rm sym}_{\alpha,s}}^2
   &\le C\varepsilon^{5-2\alpha}(1+|\log\varepsilon|)^s.
 \end{aligned}
 \tag{MA-PI.271}
\]
Both tend to zero.  The annular integral correction still
has size \(O(\varepsilon^3)\).  The Fourier derivative
estimate (MA-PI.224) is unchanged, so the cutoff error
tends to zero in Schwartz space and then under \(V\)
in the given norm.  No whole-space cutoff bound is used.

Choose an open interval \(I\subset(0,\infty)\) on which
\(0<\rho_Y<1\).  Three nonzero even smooth bumps with
pairwise disjoint supports inside \(I\cup(-I)\) have a
nontrivial linear combination \(u\) satisfying
\[
 \int u(\xi)\,d\xi=0,\qquad
 \int\rho_Y(|\xi|)u(\xi)\,d\xi=0,\qquad
       (\rho_Y^2-\rho_Y)u\ne0.
 \tag{MA-PI.272}
\]
There are only two linear constraints; disjoint supports
ensure the resulting nonzero combination is not zero,
and the transition factor is nonzero on its support.
For \(f=\mathcal F^{-1}u\), both \(f\) and \(B_Yf\)
belong to \(\mathscr X_0\): their origin values are the
two integrals in (MA-PI.272), and their integrals vanish
because \(u\) is supported away from zero.

Apply the extended full equation to \(f\) and to \(B_Yf\).
Idempotence of the actual global projector gives
\[
 R V(B_Y^2-B_Y)f
      =P_{Y,\mathscr H}^{2}RVf-P_{Y,\mathscr H}RVf=0.
 \tag{MA-PI.273}
\]
Injectivity of \(R\) and \(V\) contradicts the last
condition of (MA-PI.272).  Thus the full core equation
fails at this profile and level, as asserted.
\end{proof}

\paragraph{An explicit nonzero flat outside input.}
The outside part of (MA-PI.255) also has nonzero genuine
inputs.  Define
\[
 u(t)=
 \begin{cases}
 t^{-1}e^{-(\log t)^2}\sin(2\pi\log t),&t>0,\\
 0,&t\le0.
 \end{cases}
 \tag{MA-PI.274}
\]
Every derivative on the positive ray is a finite sum of
\(t^{-j-1}\) times a polynomial in \(\log t\), the same
Gaussian, and a sine or cosine.  This proves that the
extension is Schwartz and flat at zero.  Substitution
\(x=\log t\) and the Gaussian integral give
\[
 \int_0^\infty t^n u(t)\,dt
   =\sqrt\pi e^{n^2/4-\pi^2}\sin(n\pi)=0,
                              \qquad n=0,1,2,\ldots.
 \tag{MA-PI.275}
\]
For a fixed \(Y\), choose \(A_0>2Y\), and set
\(v(\xi)=u(|\xi|-A_0)\).  This is a nonzero even
Schwartz function supported in \(\{|\xi|\ge A_0\}\).
Its polynomial moments all vanish, by parity and the
binomial expansion of \((t+A_0)^n\) in (MA-PI.275).
Thus \(f=\mathcal F^{-1}v\) is a nonzero even Schwartz
function with every origin derivative zero.  Its integral
is \(v(0)=0\).  Hence \(f\in\mathscr X_{\rm fl}\),
\(B_Yf=0\), and injectivity proves \(RVf\ne0\).
This source is Schwartz and flat; no compact-support
assertion is made about it.

\paragraph{Source, metric, and logical boundaries.}
The local chart, exact factor two and Haar normalization
are those of MA-PI-010, from Masullo pp.~99--100, 434,
515 and~556,560--562.  The local smooth space on p.~515
is not itself identified with its Hilbert completion.
The one-place PW source and its actual projection are
the exact definitions cited before (MA-PI.269).
No extra unit-shell transition set, finite Euler deletion,
or shifted physical level is inserted.

The local functional equation and real table are Tate
(1950), Theorem~2.4.1 and \S2.5, thesis pp.~(2.11) and
(2.14), scan PDF pp.~18 and~21.  The actual Fourier
convention in (MA-PI.250) and the even sector fix the
normalization used here.  The statement/table interface
and its scalar conversion were checked; no new full
audit of Tate's thesis is claimed.

For the printed family the actual unitary to the source
space is \(U_bu(\nu(e^x)k)=m^{-1/2}e^{bx}u(x)\), with
the retained compact Haar mass \(m\) and \(b=A-1/2\).
It transports the same norm, closed PW range and
orthogonal projector.  For the symmetric family the
norm in (MA-PI.251) is explicitly the corrected metric.
The proof does not move an orthogonal projector by an
arbitrary equivalent norm.  No endpoint \(|b|=1/2\)
or \(\alpha=1/2\) extension is asserted.

The output of (MA-PI.253) is in the actual closed ideal;
it is not asserted to be in the full smooth ideal for
every completed input.  The first plateau in (MA-PI.255)
is vacuous.  Thus the example neither contradicts the
full-Schwartz two-sided plateau theorems nor provides
the full intertwining equation.  Step~6 positively
separates those scopes in every norm under consideration.
It supplies no new trace comparison, inverse, nuclear
endomorphism trace, or arithmetic vanishing in MA-FT-002.
This newly formulated counterexample is not an old-row
movement witness.

In the inherited cote~1 BR-TR-01/02 placement,
scans079--080 and111--112, this is an actual bounded
map between specified completed objects, with a precisely
limited compatibility.  The failure of the stronger
idempotent relation is an operator-domain statement,
not a consequence of a finite-length \(K_0\) class.
Norm-half centering remains analytic, not an integral
Tate twist; no additional weight, purity or positivity
claim is introduced.
